Eine feste Menge eines idealen Gases durchläuft einen Kreisprozess. Vom Zustand A ($p_A = 2\,p_0$, $V_A = 2\,V_0$) wird es isochor erwärmt, bis sich der Druck verdreifacht hat (Zustand B). Danach dehnt es sich isotherm aus, bis das Volumen dreimal so gross ist (Zustand C), und wird schliesslich isobar zurück in den Zustand A komprimiert. Die Achsen sind in Einheiten von $p_0$, $V_0$ und $T_0 = \frac{p_0V_0}{nR}$ beschriftet; damit ist $T/T_0 = (p/p_0)\cdot(V/V_0)$.

Eine Schülerin zeichnet dazu das folgende $p(V)$-Diagramm:

\begin{center}
\begin{tikzpicture}[x=0.75cm, y=0.75cm, >=Stealth, font=\small]
 \draw[black!12] (0,0) grid[xstep=1, ystep=1] (7.001,7.001);
 \draw[->] (0,0) -- (7.6,0) node[right] {$V$ in $V_0$};
 \draw[->] (0,0) -- (0,7.6) node[above] {$p$ in $p_0$};
 \node[below left] at (0,0) {\small 0};
 \draw (1,0.08) -- (1,-0.08) node[below] {\small 1};
 \draw (2,0.08) -- (2,-0.08) node[below] {\small 2};
 \draw (3,0.08) -- (3,-0.08) node[below] {\small 3};
 \draw (4,0.08) -- (4,-0.08) node[below] {\small 4};
 \draw (5,0.08) -- (5,-0.08) node[below] {\small 5};
 \draw (6,0.08) -- (6,-0.08) node[below] {\small 6};
 \draw (7,0.08) -- (7,-0.08) node[below] {\small 7};
 \draw (0.08,1) -- (-0.08,1) node[left] {\small 1};
 \draw (0.08,2) -- (-0.08,2) node[left] {\small 2};
 \draw (0.08,3) -- (-0.08,3) node[left] {\small 3};
 \draw (0.08,4) -- (-0.08,4) node[left] {\small 4};
 \draw (0.08,5) -- (-0.08,5) node[left] {\small 5};
 \draw (0.08,6) -- (-0.08,6) node[left] {\small 6};
 \draw (0.08,7) -- (-0.08,7) node[left] {\small 7};
 \draw[very thick, blue] (2,2) -- (2,6);
 \draw[-{Stealth[length=2.6mm]}, blue] (2,4.19) -- (2,4.21);
 \draw[very thick, blue] (2,6) -- (6,2);
 \draw[-{Stealth[length=2.6mm]}, blue] (4.193,3.807) -- (4.207,3.793);
 \draw[very thick, blue] (6,2) -- (2,2);
 \draw[-{Stealth[length=2.6mm]}, blue] (3.81,2) -- (3.79,2);
 \fill[black] (2,2) circle (2pt) node[below left] {$A$};
 \fill[black] (2,6) circle (2pt) node[above right] {$B$};
 \fill[black] (6,2) circle (2pt) node[above right] {$C$};
\end{tikzpicture}
\end{center}

\begin{abcliste}
\abc Bestimmen Sie Druck, Volumen und Temperatur in den Zuständen A, B und C. Stimmen die drei Eckpunkte der Skizze?
\abc Welcher Schritt ist falsch gezeichnet? Zeigen Sie mit dem Mittelpunkt $M$ dieser Linie, dass sie keine Isotherme sein kann.
\abc Skizzieren Sie den Schritt richtig. Wo liegt er bei $V = 3\,V_0$, $4\,V_0$ und $5\,V_0$?
\end{abcliste}
